\section*{Chapter \ref{ch:m1:basis-roll-and-delivery} --- Basis, Roll and Delivery}

\begin{solution}{exo:m1:basis-roll-and-delivery:1}
$6\,000 \times (1 + 0.042 \times 91/360) = 6\,063.70$; \omterm{def:m1:basis-roll-and-delivery:basis}{basis} 63.70 points.
\end{solution}

\begin{solution}{exo:m1:basis-roll-and-delivery:2}
$6.0 \times (1 + 0.042 \times 64/360) + 9.5 \times (1 + 0.042 \times 32/360) + 5.5
\times (1 + 0.042 \times 8/360) = 6.04 + 9.54 + 5.51 = 21.09$. \omterm{def:m1:basis-roll-and-delivery:carry}{Fair value}
$6\,063.70 - 21.09 = 6\,042.61$.
\end{solution}

\begin{solution}{exo:m1:basis-roll-and-delivery:3}
\omterm{def:m1:basis-roll-and-delivery:contango}{Contango}; a long \omterm{def:m1:basis-roll-and-delivery:roll}{rolls} from 70.00 into 70.84 each month and, if nothing
moves, sees 70.84 become 70.00: $(70/70.84 - 1) \times 12 = -14.2\%$ a year.
\omterm{def:m1:basis-roll-and-delivery:contango}{Backwardation}; $(70/69.30 - 1) \times 12 = +12.1\%$ a year.
\end{solution}

\begin{solution}{exo:m1:basis-roll-and-delivery:4}
The future is $6\,050.50 - 6\,042.61 = 7.89$ points rich, beyond the upper
bound of $+6.47$. Sell the future at 6\,050.50; buy the basket, \$300\,000 per
contract, in index proportions; finance it at the benchmark plus 15 basis
points; collect and reinvest the dividends; on expiry morning sell every
stock in its opening auction. Locked in: $7.89 - 6.47 = 1.41$ points, \$70 a
contract. What can go wrong: legging (the basket of 500 stocks is not bought
in one instant); dividends cut or moved; the financing rate rising; stocks
halted at the open on expiry day; \omterm{def:m1:clearing-and-settlement:margin}{variation margin} on the short future
during a rally, to be funded while the basket's gain is unrealised.
\end{solution}

\begin{solution}{exo:m1:basis-roll-and-delivery:5}
Bisection on $F^*(r) = 6\,047.20$ gives $r = 4.50\%$: check, $6\,000 \times
0.0450 \times 91/360 = 68.25$, dividends 21.09, $F^* = 6\,047.16$, the rest
being rounding of the rate. A rise of 25 basis points adds $6\,000 \times
0.0025 \times 91/360 = 3.79$ points of interest and a negligible amount of
reinvestment: the \omterm{def:m1:basis-roll-and-delivery:carry}{fair value} rises by 3.79 points with the index unchanged.
Futures are interest-rate instruments too.
\end{solution}

\begin{solution}{exo:m1:basis-roll-and-delivery:6}
The published index at 09:30 mixes stocks that have opened with the
previous close of those that have not; it is 5\,990 only by that convention.
The quotation uses each stock's own opening print, whenever it occurs. If
the stocks that open late open low, their opening prices enter the
quotation but their contribution to the ``day's low'' of the published index
comes later, at other stocks' then-current prices: the quotation is a
portfolio of prices taken at different times and need not lie within the
range of any simultaneous index. Holders of expiring futures and options
settle on it; arbitrageurs who sell their baskets in the opening auctions
receive exactly it.
\end{solution}

\begin{solution}{exo:m1:basis-roll-and-delivery:7}
0.751 in \omterm{def:m1:basis-roll-and-delivery:contango}{contango} and 1.273 in \omterm{def:m1:basis-roll-and-delivery:contango}{backwardation} after 24 monthly \omterm{def:m1:basis-roll-and-delivery:roll}{rolls}. With a
\omterm{def:m1:basis-roll-and-delivery:contango}{contango} shrinking linearly from 1.2\% to zero: 0.861. Half the slope on
\omterm{def:m1:pnl-and-positions:realised}{average costs} roughly half the loss.
\end{solution}

\begin{solution}{exo:m1:basis-roll-and-delivery:8}
The premium is carry: three months of interest less three months of
dividends, 42.6 points at these rates, and it would be the same if everyone
expected a crash. Expectations move the index and the future together. What
the future tells you before the open is the \emph{change}: compare the
future with yesterday's \omterm{def:m1:basis-roll-and-delivery:carry}{fair value} at the close (yesterday's close plus the
fair \omterm{def:m1:basis-roll-and-delivery:basis}{basis}); that difference, not the raw premium, is the market's overnight
move.
\end{solution}

\begin{solution}{pb:m1:basis-roll-and-delivery:1}
\textbf{1.} $1\,000\,000\,000/300\,000 = 3\,333$ contracts.
\textbf{2.} 6\,042.61.
\textbf{3.} 182 days of interest, 127.40 points, less all six dividends
carried to 19 March, 42.39 points: 6\,085.01.
\textbf{4.} 42.39 points.
\textbf{5.} Interest for the 91 days between the expiries on 6\,000: 63.70.
Dividends of that quarter, with reinvestment: 21.31. $63.70 - 21.31 = 42.39$.
\textbf{6.} $4.55/6\,000/(91/360) = 30$ basis points a year.
\textbf{7.} $4.55 \times 50 \times 3\,333 = \$758\,258$.
\textbf{8.} $0.025 \times 50 \times 3\,333 = \$4\,166$.
\textbf{9.} Four \omterm{def:m1:basis-roll-and-delivery:roll}{rolls}: \$3\,049\,695, 30.5 basis points of the exposure.
\textbf{10.} The futures cost 30.5 against the fund's 7, but the fund's cash
earns the benchmark in both cases only if the fund is fully funded; futures
leave 95\% of the cash free (for a fund that can earn more than the benchmark
on it, or that needs it as collateral elsewhere), avoid \omterm{def:m1:delta-one-instruments:wht}{withholding tax}
differences, and can be shorted. For a fully funded long-only holder at these
levels the index fund is cheaper.
\textbf{11.} With the crowd: that is when the spread book is deepest and the
tick is a small cost; the richness is whatever it is. Earlier: a holder who
believes the richness rises into the \omterm{def:m1:basis-roll-and-delivery:roll}{roll} week (as dealers' balance sheets
fill) pays less by going first, at the price of a thinner book.
\textbf{12.} The March contract spans the year-end, when banks shrink their
balance sheets for reporting dates and financing is scarcest; the December
\omterm{def:m1:basis-roll-and-delivery:roll}{roll} prices that.
\textbf{13.} The \omterm{def:m1:basis-roll-and-delivery:basis}{basis} itself, in index points against a known close, hence
the \omterm{ex:m1:basis-roll-and-delivery:43}{implied financing rate} to expiry without having to synchronise a futures
price with an index that is published with a delay.
\textbf{14.} It can, because it holds or can finance the basket: it sells
the future rich and earns benchmark plus 20. It would, because 20 basis
points on a balance-sheet position is its business; it will not at year-end
if the capital cost exceeds that.
\textbf{15.} The slope is not an interest rate but storage and convenience
yield, it can be many percent a year, it changes sign, and the choice of
which contract to hold along the curve becomes a strategy in itself.
\textbf{16.} The E-mini settles in cash against the special opening
quotation: the fund receives or pays a last \omterm{def:m1:clearing-and-settlement:margin}{variation margin} and is out of
the market, unhedged, from that morning. WTI is delivered: a long who has not
closed is assigned barrels at Cushing and must have arranged pipeline or
storage; brokers close such positions by force before it comes to that.
\textbf{17.} It is the market price of dealer balance sheet for equity
financing, the same number that appears in swap spreads, in the cost of
leveraged strategies and in the returns of index arbitrage.
\textbf{18.} Whoever is short the future and long the stock: index
arbitrageurs, dealers, and funded investors who replace their shares with
futures when the \omterm{def:m1:basis-roll-and-delivery:roll}{roll} is cheap and the reverse when it is rich.
\textbf{19.} \textbf{30.5 basis points a year.}
\textbf{20.} Because at expiry it pays the index, so anyone who can hold the
basket until then can collect any difference.
\end{solution}

\begin{solution}{iq:m1:basis-roll-and-delivery:1}
The buyer of the future does not pay for the shares and earns interest on
the cash, but receives no dividends; \omterm{def:m1:basis-roll-and-delivery:carry}{fair value} is spot plus interest minus
dividends. It trades at a discount to spot when the dividend yield to expiry
exceeds the interest rate, as in low-rate years or in markets with a heavy
dividend season before expiry, and, relative to \omterm{def:m1:basis-roll-and-delivery:carry}{fair value}, when the basket
is costly to short.
\iqlookfor{carry, not expectations; a concrete case of discount.}
\end{solution}

\begin{solution}{iq:m1:basis-roll-and-delivery:2}
\omterm{def:m1:basis-roll-and-delivery:contango}{Contango}: later expiries above nearer ones; \omterm{def:m1:basis-roll-and-delivery:contango}{backwardation}: below. Neither
is a forecast. For a financial asset the slope is rate minus yield; for a
commodity it adds storage and subtracts convenience yield. A curve in
\omterm{def:m1:basis-roll-and-delivery:contango}{contango} says that carrying the commodity is costly and inventories are
ample; a long investor pays the slope as negative \omterm{def:m1:basis-roll-and-delivery:roll}{roll} yield.
\iqlookfor{slope as carry, and the roll-yield consequence.}
\end{solution}

\begin{solution}{iq:m1:basis-roll-and-delivery:3}
Compute \omterm{def:m1:basis-roll-and-delivery:carry}{fair value} from the rate and dated dividends. When the future is
above \omterm{def:m1:basis-roll-and-delivery:carry}{fair value} by more than my costs, sell futures and buy the basket
through a program trade, financed. Carry: collect dividends, pay financing,
pay or receive \omterm{def:m1:clearing-and-settlement:margin}{variation margin}. Exit either when the mispricing reverses
(sell the basket, buy the future) or at expiry, by selling each stock in its
opening auction on expiry morning, which realises the special opening
quotation exactly. Risks: legging, dividends, financing, halts, margin cash.
\iqlookfor{the opening-auction exit and the margin cash on the short future.}
\end{solution}

\begin{solution}{iq:m1:basis-roll-and-delivery:4}
(i) Unadjusted front month: true traded prices, with a jump at each \omterm{def:m1:basis-roll-and-delivery:roll}{roll}; use
for levels on a given day, never for returns. (ii) Difference-adjusted
(back-adjust by the spread at each \omterm{def:m1:basis-roll-and-delivery:roll}{roll}): preserves point P\&L, distorts
percentage returns and can go negative. (iii) Ratio-adjusted: preserves
percentage returns, distorts levels. (iv) Best: compute returns from
same-contract prices and chain them, keeping the contract identity; derive
any level series from that. The \omterm{def:m1:basis-roll-and-delivery:roll}{roll} date rule (volume or open-interest
cross, or fixed days) must be explicit and must be one the strategy could
have followed.
\iqlookfor{returns from same-contract prices; awareness that ratio and
difference adjustments answer different questions.}
\end{solution}

\begin{solution}{iq:m1:basis-roll-and-delivery:5}
The contract is \omterm{def:m1:basis-roll-and-delivery:delivery}{physically delivered} at a landlocked hub. A day before
expiry, longs who could not take delivery had to sell, storage at the hub was
committed, and those able to take oil demanded payment to do so; with much
of the day's volume tied to the settlement through trade-at-settlement
orders, the settlement window had sellers and almost no buyers. A
\omterm{def:m1:basis-roll-and-delivery:delivery}{cash-settled} index future cannot do this: its final price is the index,
which is non-negative, and nobody is ever obliged to store anything.
\iqlookfor{delivery and storage as the mechanism; \omterm{def:m1:basis-roll-and-delivery:delivery}{cash settlement} as the
difference.}
\end{solution}

\begin{solution}{iq:m1:basis-roll-and-delivery:6}
Sell the \omterm{def:m1:basis-roll-and-delivery:roll}{roll}: buy December, sell March, against a long basket or, more
simply, be short March futures against shares held from December's expiry
onward, earning benchmark plus 60 over the year-end. Equivalent forms: sell
BTIC, lend through a \omterm{def:m1:financing:trs}{total return swap}. Limits: it consumes balance sheet
and capital over the reporting date, which is exactly why it is rich; the
position is long \$X of stock financed over year-end; \omterm{def:m1:pnl-and-positions:mtm}{mark-to-market} of the
spread if richness widens further; dividend risk; and for a fund, the
variation-margin cash on the short future. Size to the balance sheet one can
commit across the turn, not to the spread.
\iqlookfor{recognising that the richness is the price of the constraint
that limits the trade.}
\end{solution}
